The results are presented as an evidence chain: independent holdout accuracy, exact replay of selected decisions, operating consequences, objective trade-offs, and primary search time. `N` denotes the nominal influent and `S1`--`S10` denote the post-selection influent scenarios. Unless stated otherwise, selected-decision values are exact mechanistic replays rather than optimizer-native surrogate predictions.

\subsection{Independent holdout validation of the projected surrogate}

Figure~\ref{fig:holdout-overview} summarizes independent holdout accuracy before and after physical projection. Projection lowers aggregate coordinate-normalized RMSE from 0.0469 to 0.0431 and coordinate-normalized MAE from 0.0304 to 0.0257, while raising mean location $R^2$ from 0.925 to 0.938. The projected response is therefore the appropriate object for the subsequent optimization assessment.

\begin{figure}[htbp]
\centering
\includegraphics[width=\textwidth]{../../results/article_v3/article_full_10000_001/report/figures/q01_holdout_accuracy_overview.png}
\caption{Independent holdout accuracy before and after physical projection. Rows report aggregate, location-level, and water-quality-composite metrics. nRMSE and nMAE are coordinate-normalized and lower values are better; mean location $R^2$ is higher values are better.}
\label{fig:holdout-overview}
\end{figure}

The location-by-composite decomposition in Figure~\ref{fig:holdout-location} places the aggregate improvement in the connected plant. Projected nRMSE is lower in 30 of 32 location--composite cells. Retaining the two exceptions and remaining small changes makes the accuracy statement auditable at the mixer, biological stages, and clarifier outlets.

\begin{figure}[htbp]
\centering
\includegraphics[width=\textwidth]{../../results/article_v3/article_full_10000_001/report/figures/q02_holdout_accuracy_by_location.png}
\caption{Location-by-composite holdout nRMSE before and after physical projection. The right panel reports percentage change; negative values denote an improvement.}
\label{fig:holdout-location}
\end{figure}

Figure~\ref{fig:holdout-parity} supplies the complementary observation-level view. It pools the eight locations for each composite while retaining an explicit median marker for every location; hexagon colour reports local point density, not a third prediction method. The agreement follows the one-to-one line across the observed ranges. Projected nRMSE is 0.0345 for TN, 0.0335 for TP, 0.0459 for COD, and 0.0548 for TSS. Together, Figures~\ref{fig:holdout-overview}--\ref{fig:holdout-parity} separate independent prediction validation from the post-selection results.

\begin{figure}[htbp]
\centering
\includegraphics[width=\textwidth]{../../results/article_v3/article_full_10000_001/report/figures/q03_holdout_parity_all_locations.png}
\caption{Projected Extended-ICSOR parity across all holdout observations and system locations. Hexagons show location-observation counts; explicitly labelled markers identify location medians.}
\label{fig:holdout-parity}
\end{figure}

\FloatBarrier
\subsection{Exact replay of selected Extended-ICSOR decisions}

Independent holdout accuracy does not alone establish fidelity at an optimized decision. Figure~\ref{fig:selected-parity} compares the Extended-ICSOR values used during selection with exact replay of the same nominal and scenario decisions. Concentration and removal parity appear together because a small concentration difference can correspond to a different removal percentage when the influent changes. The median absolute concentration percentage error is 12.7\%, and the median absolute removal difference is 3.26 percentage points. These are post-selection descriptive errors, not holdout errors; exact replay is consequently retained as the common basis for route comparison.

\begin{figure}[htbp]
\centering
\includegraphics[width=\textwidth]{../../results/article_v3/article_full_10000_001/report/figures/q04_effluent_and_removal_parity.png}
\caption{Extended-ICSOR quantities at selected decisions compared with exact mechanistic replay. The upper row reports effluent concentrations and the lower row reports removals for the nominal and all influent scenarios.}
\label{fig:selected-parity}
\end{figure}

\FloatBarrier
\subsection{Operating outcomes and treatment-train behavior}

Figure~\ref{fig:effluent-controls} places exact replayed effluent outcomes beside the controls that generated them: HRT, aeration, recycle, return sludge, and wasting. The nominal decision is displayed with every influent scenario, which separates the operating baseline from sensitivity to the tested influent set. The median absolute route difference in exact effluent concentrations is 21.4\%; this descriptive contrast is not a reliability estimate beyond the evaluated scenarios.

\begin{figure}[htbp]
\centering
\includegraphics[width=\textwidth]{../../results/article_v3/article_full_10000_001/report/figures/q05_effluent_and_operating_values.png}
\caption{Exact-replay effluent concentrations and selected operating controls for Extended ICSOR and Smooth NLP. `N` is the nominal influent; `S1`--`S10` are influent scenarios.}
\label{fig:effluent-controls}
\end{figure}

Figure~\ref{fig:treatment-profiles} connects those decisions to process behavior. Each composite has a common vertical scale for the two routes, so within-composite route differences are visible without confusing them with the different magnitudes of COD, TN, TP, and TSS. The profiles follow each scenario from influent through the mixer and R1--R5 to clarifier effluent. Here R1--R5 identify biological stages, whereas S1--S10 identify influent scenarios.

\begin{figure}[htbp]
\centering
\includegraphics[width=\textwidth]{../../results/article_v3/article_full_10000_001/report/figures/q06_treatment_train_profiles.png}
\caption{Exact-replay COD, TN, TP, and TSS profiles through the treatment train for the nominal influent and all influent scenarios. For each composite, the Extended-ICSOR and Smooth-NLP panels share logarithmic y-axis limits.}
\label{fig:treatment-profiles}
\end{figure}

\FloatBarrier
\subsection{Objective trade-offs and primary search time}

Figure~\ref{fig:objective-tradeoffs} resolves the exact replay objective into total objective, normalized water-quality contribution, and weighted resource contributions. In the ten influent scenarios, Smooth NLP has a lower exact total objective in all ten and a lower normalized water-quality component in nine. These paired results describe the evaluated candidates; they do not confer a global optimality certificate on either route.

\begin{figure}[htbp]
\centering
\includegraphics[width=\textwidth]{../../results/article_v3/article_full_10000_001/report/figures/q07_objective_quality_economic.png}
\caption{Exact-replay objective comparison for the nominal influent and all influent scenarios. Panels show total objective, normalized water-quality component, and weighted resource contributions.}
\label{fig:objective-tradeoffs}
\end{figure}

Figure~\ref{fig:optimization-time} puts that objective comparison in computational context. Extended ICSOR is faster in every influent scenario, with mean primary-search times of 1.73~s for Extended ICSOR and 5.24~s for Smooth NLP. The logarithmic axis shows scenario-to-scenario spread; data generation, fitting, certification, recovery, and exact replay are excluded by the definition of Time. This is a comparison of the stated search procedures on the common computing allocation, not total workflow time.

\begin{figure}[htbp]
\centering
\includegraphics[width=\textwidth]{../../results/article_v3/article_full_10000_001/report/figures/q08_optimization_time.png}
\caption{Primary optimization time for Extended ICSOR and Smooth NLP. The nominal influent and all influent scenarios are shown; the y-axis is logarithmic.}
\label{fig:optimization-time}
\end{figure}

\FloatBarrier
